\begin{definition}[Restoring the dimension multiplicities]%
    \label{def:peps_multiplicity_restored_representation}
    \lean{TNLean.PEPS.multiplicityRestoredRepresentation}
    \leanok
    Given matrix representations $D_i:G\to M_{d_i}(\mathbb C)$, define
    \begin{align}
        L(g)=\bigoplus_iD_i(g)\otimes I_{d_i}.
    \end{align}
    This is a representation. If the $D_i$ are precisely the inequivalent
    irreducible representations, these are the dimension multiplicities
    in the regular representation used in \cite[Section~7]{Schuch2010PEPS}.
    The definition itself does not assert that every irreducible occurs.
\end{definition}

\begin{theorem}[Weighted block identities and the actual torus state]%
    \label{thm:peps_torus_weighted_blocks_state}
    \lean{TNLean.PEPS.torusBondRegrouping_dressedAveragingSite_of_weightedBlocks}
    \leanok
    \uses{thm:peps_coherent_multiplicity_transport,
        def:peps_multiplicity_restored_representation,
        thm:peps_torus_regrouped_averaging_bonds,
        lem:peps_matching_sector_bond_inclusion,
        thm:peps_full_multiplicity_bond_map}
    Let $d_i>0$, let $U$ and $D_i$ be matrix representations of a finite group,
    and let $W$ commute with every $U(g)$. Suppose that the single-bond matrices satisfy
    \begin{align}
        W^2U(g)=\bigoplus_i\sqrt{d_i}D_i(g) \qquad(g\in G).
    \end{align}
    The actual dressed averaging-site torus state $H$ then satisfies
    \begin{align}
        \mathcal B H               & \in\operatorname{range}(E^{\otimes 2N}), &
        F^{\otimes 2N}\mathcal B H & =\mathcal B K_L.
    \end{align}
    Here $\mathcal B$ regroups physical legs into bonds, $E$ includes matching sectors,
    $F$ restores the dimension multiplicities, and $K_L$ is the actual averaging-site
    state of $L(g)=\bigoplus_iD_i(g)\otimes I_{d_i}$.
    This auxiliary assertion assumes a matrix identity, rather than an equality
    of tensor-network states; it applies to both constructions below.
\end{theorem}

\begin{proof}\leanok
    Its coherent expansion has bond factors $W^2U(r_e(q))$.
    Matching-sector inclusion reconstructs each weighted block
    matrix, and multiplicity restoration sends it to $L(r_e(q))$.
    Apply these identities to every term of the finite vertex-label sum.
    Both contractions have the same normalization $|G|^{-N}$.
\end{proof}

\begin{theorem}[Multiplicity restoration of the actual torus state]%
    \label{thm:peps_torus_multiplicity_bond_state}
    \lean{TNLean.PEPS.physicalProductMap_sum_prod,
        TNLean.PEPS.torusBondRelativeElement,
        TNLean.PEPS.torusBondRegrouping_averagingSite_coherent,
        TNLean.PEPS.torusBondRegrouping_dressedAveragingSite_coherent,
        TNLean.PEPS.exists_torusMultiplicityBondState_of_orthogonalSectors}
    \leanok
    \uses{def:peps_multiplicity_restored_representation,
        thm:peps_theta_bond_orthonormal_coordinates,
        thm:peps_torus_weighted_blocks_state,
        thm:peps_torus_regrouped_averaging_bonds,
        lem:peps_matching_sector_bond_inclusion,
        thm:peps_full_multiplicity_bond_map}
    Let $\rho$ act on a finite-dimensional complex Hilbert space with a
    supplied orthogonal internal decomposition into irreducible sectors
    $S_i$, each of character multiplicity one. Choose an orthonormal basis
    of each $S_i$, of size $d_i$, and let $D_i(g)$ be its representation matrix.
    There is a global orthonormal basis $c$ collecting these sector coordinates.
    Put $U(g)=[\rho(g)]_c$ and $W=[\Theta]_c$.

    Let $H_c$ be the actual torus state obtained by dressing the normalized
    averaging site of $U$ with $W$ on every virtual leg. Let $K_L$ be the
    actual normalized averaging-site torus state of the repeated representation
    $L$. Write $\mathcal B$ for physical bond regrouping, $E$ for matching-sector
    inclusion, and $T$ for the full endpoint multiplicity-restoring map.
    On a torus with positive circumferences and $N$ vertices,
    \begin{align}
        \mathcal B H_c               & \in\operatorname{range}(E^{\otimes 2N}), &
        T^{\otimes 2N}\mathcal B H_c & =\mathcal B K_L.
    \end{align}
    The support and equality are conclusions about the actual contractions.
    This auxiliary assertion uses a supplied decomposition; it does not
    establish the existence of the smallest semi-regular representation or
    its Fourier identification with the regular representation. These steps
    are separate; see \cite{gap:rmp_peps_quantum_double_g_isometry}.

    For the oriented torus bond set $B=V_{\rm tor}\times\{0,1\}$, write
    \begin{align}
        r_{(v,0)}(q) & =q_{v+\mathrm{east}}q_v^{-1},  &
        r_{(v,1)}(q) & =q_vq_{v+\mathrm{north}}^{-1}.
    \end{align}
    The regrouped averaging-site state has the actual coherent expression
    $|G|^{-N}\sum_q\bigotimes_{e\in B}U(r_e(q))$.
    A commuting dressing $W$ replaces each bond factor by $W^2U(r_e(q))$.

    More generally, for a finite bond set $B$, any rectangular physical matrix
    $F$ acts on arbitrary coherent sums by
    \begin{align}
        F^{\otimes B}\left(\sum_q w_q\bigotimes_{e\in B}x_{e,q}\right)
        =\sum_q w_q\bigotimes_{e\in B}Fx_{e,q}.
    \end{align}
\end{theorem}

\begin{proof}\leanok
    Expand the actual dressed site contraction after physical bond regrouping.
    Each relative group element contributes a matrix
    $W^2U(g)=\bigoplus_i\sqrt{d_i}D_i(g)$.
    Matching-sector inclusion reconstructs this matrix from its diagonal
    blocks, proving the product support assertion. On every such bond vector,
    multiplicity restoration gives $\bigoplus_iD_i(g)\otimes I_{d_i}=L(g)$.
    The coherent-sum identity follows by interchanging the finite sums and
    factoring the independent sums over bond labels. Apply it to the vertex
    group-label expansion of the actual torus state. The result is precisely
    the regrouped averaging-site contraction of $L$, with the same normalization
    $|G|^{-N}$ on both sides.
\end{proof}
